\chapter{Artificial Intelligence: The Future Is Here}
\label{ch:ai}

By now you have learned to explore data, reason about probability,
test hypotheses, estimate unknown quantities, collect data responsibly,
model relationships between variables, and work through a full case
study from raw data to inference. Along the way, a new kind of
collaborator has become part of almost every data analyst's toolkit:
artificial intelligence, or AI. Large language models and related tools
can now write code, suggest analyses, produce visualizations, and
explain statistical ideas in plain language. Used well, AI can make you
faster and help you learn; used carelessly, it can quietly introduce
errors, invalid conclusions, or hallucinated results into your work. In
this chapter, we step back from R code and discuss how to work with AI
as a data analyst: what it tends to do well, where it tends to go
wrong, and how the judgment you have built up over the previous
chapters---about representative samples, bias, appropriate methods,
and the limits of observational data---is exactly what keeps AI useful
rather than dangerous.

This chapter contains no R code, and that is deliberate, for a few
reasons. Not everyone has access to the same AI tools: some require a
paid account, while R itself is free and available to everyone.
Different AI tools behave differently from one another, and even the
same tool will rarely produce identical code, text, or analysis twice.
Most importantly, specific tools and interfaces change quickly, but
the principles in this chapter---what AI is good at, what it is bad
at, and how to keep yourself in charge of the analysis---should stay
relevant regardless of which particular AI you use, now or in a few
years.

\section{What AI Can Do for You}

Used as a collaborator rather than an oracle, AI can meaningfully
speed up almost every stage of a data analysis project.

\begin{itemize}
\item \textbf{Clean, documented code.} AI tools are often good at
  producing code that is well organized and commented, following
  common style conventions. This can be a useful model to learn from,
  especially early on, and can save time on boilerplate code you would
  otherwise write by hand.
\item \textbf{Automating data acquisition.} AI can help write scripts
  to download, parse, and reshape data from application programming
  interfaces (APIs), websites, or messy file formats, tasks that are
  often tedious and repetitive.
\item \textbf{Debugging, testing, and validation.} AI is often good at
  reading error messages, suggesting likely causes, and proposing
  fixes. It can also help write small test cases to check that a
  function behaves as expected on edge cases you might not have
  thought of.
\item \textbf{Improving writing.} AI can help you explain a result
  more clearly, tighten a paragraph, or check that your statistical
  language is precise (for example, distinguishing an association from
  a causal effect).
\item \textbf{Producing visualizations.} AI can quickly generate a
  first draft of a plot, including plots you have not made before, and
  suggest alternative ways to display the same data.
\item \textbf{Borrowing skills from experienced people.} Because AI is
  trained on an enormous amount of code and writing from many
  practitioners, asking it for a second opinion is a little like
  asking a large group of experienced analysts how they would approach
  a problem.
\end{itemize}

\section{Where AI Can Lead You Astray}

The same qualities that make AI powerful also make it risky to use
without care.

\begin{itemize}
\item \textbf{Hallucination.}\index{hallucination} AI tools sometimes
  invent function names, arguments, citations, or numerical results
  that look plausible but do not exist or are simply wrong. A
  hallucinated result can be stated just as confidently as a correct
  one, so it is not always obvious which is which.
\item \textbf{Reproducibility.}\index{reproducibility} An AI
  assistant is not a substitute for version control or a documented
  workflow. If you regenerate an analysis with a different prompt, a
  different tool, or even the same tool on a different day, you may
  not get the same code or the same answer. You, not the AI, must be
  the one who saves the final code, tracks changes, and can explain
  and rerun the whole analysis later.
\item \textbf{False confidence.} AI tends to state its answers in a
  fluent, assured tone whether or not the underlying reasoning is
  correct. This is sometimes called \emph{automation
  bias}\index{bias!automation bias}: the tendency to trust a
  computer-generated answer more than it deserves, simply because it
  is delivered so confidently. Nothing about the tone of a response
  tells you how reliable it is.
\item \textbf{Data leakage and privacy.}\index{data leakage} Pasting a
  dataset, or even a small snippet of one, into a general-purpose AI
  tool may send that data outside your organization or beyond the
  scope of what your data-use agreement allows. This is a particular
  concern for medical records, student data, financial records, or any
  data covered by a privacy rule, an institutional review board, or a
  data-sharing agreement.
\end{itemize}

\section{Putting This Book's Skills to Work with AI}

The most effective way to use AI in a data analysis project is not to
hand it a dataset and ask for ``the answer.'' It is to bring the
skills from earlier chapters to the conversation, so that AI's speed
is guided by your judgment rather than replacing it.

\subsection{Teaching AI Your Data-Cleaning Checklist}

Chapter~\ref{ch:data} introduced the idea of a representative sample
and the difference between a controlled experiment and an
observational study, and Chapter~\ref{ch:EDA} introduced the plots and
tables (boxplots, histograms, cross-tabulations, and so on) that help
you get to know a new dataset before modeling it. Rather than asking
an AI tool to ``clean this data'' with no further guidance, give it the
specific steps you would perform yourself: check for missing values
and understand why they are missing, look for a plausible source of
selection or survival bias, produce the diagnostic plots that make
sense for the variable types involved, and flag any values that look
like data-entry errors rather than removing them outright. An AI tool
that is told exactly what to check is far more useful than one that is
asked to clean data with no criteria at all.

\subsection{Asking AI to Watch for Bias}

Chapter~\ref{ch:data} also cataloged a number of specific biases that
can silently distort a conclusion: Berkson's bias, survival bias, the
base-rate fallacy, omitted variables, observation time interval bias,
confirmation bias, availability bias, arbitrary rejection of
``outliers,'' self-selection (volunteer) bias, nonresponse bias, and
the Hawthorne effect. Rather than trusting yourself, or an AI tool, to
remember all of these while under the pressure of a deadline, give the
AI this list directly and ask it to point out which of these biases
could plausibly affect your specific dataset and research question,
and why. This turns a vague worry (``is something wrong with this
data?'') into a concrete checklist that both you and the AI can work
through together.

\subsection{Choosing and Explaining a Method}

Chapter~\ref{ch:hypothesis} showed that the right hypothesis test
depends on the structure of the data (for example, choosing a
rank-based test over a t-test when the data have heavy tails), and
Chapter~\ref{ch:regression} showed that linear regression is not
always the right model for the outcome at hand. When you ask an AI
tool to recommend a statistical method, do not simply accept the
recommendation: ask it to explain why that method fits your data and
your question, and compare that explanation to what you know about the
method's assumptions. An explanation you can follow and check is far
more valuable than a method name alone.

It is just as important to ask AI to flag a limitation you have seen
throughout this book: association is not causation, and this matters
most when the data did not come from a designed experiment. The NYC
crash case study in Chapter~\ref{ch:nyc} used a chi-squared test, a
binomial confidence interval, and a logistic regression model to
describe patterns in crash severity, but was careful to note that
these results describe an observational dataset and should not be read
as causal effects. Ask AI to add exactly this kind of caveat to its own
output, and to say explicitly whether your data came from a randomized,
controlled experiment or an observational study, since the strength of
the conclusions you are entitled to draw depends on that distinction.

\section{General Principles for Working with AI}

The tools available today will not be the tools available in a few
years, but a few working habits should remain useful regardless of
which AI you use.

\begin{itemize}
\item Treat AI as a fast, well-read, but overconfident collaborator,
  not as a source of ground truth.
\item Verify anything that matters: rerun the code yourself, check a
  cited number against its source, and confirm that a function or
  argument actually exists before relying on it.
\item Keep yourself in control of reproducibility: save the code you
  actually used, in a file or a version-control system, rather than
  relying on being able to regenerate the same output from AI later.
\item Never share data covered by a privacy rule, an institutional
  agreement, or a confidentiality obligation with a general-purpose AI
  tool unless you know exactly where that data goes.
\item Bring your own statistical judgment to every step: what makes a
  sample representative, what could bias a conclusion, and what a
  method assumes are things you should be checking, not outsourcing.
\end{itemize}

\section{Summary}

In this chapter we discussed how to work with AI as a data analyst,
without writing any R code. We surveyed AI's real strengths: producing
clean and documented code, automating data acquisition, helping with
debugging and testing, improving writing, generating visualizations,
and giving access to patterns learned from many experienced
practitioners. We also discussed its real risks: hallucinated results,
unclear reproducibility, false confidence, and the danger of data
leakage when private or sensitive data is shared with a general-purpose
tool. Finally, we connected AI back to the specific skills built in
this book: giving AI a concrete data-cleaning checklist and a list of
named biases to watch for, asking it to justify its choice of
statistical method, and requiring it to flag when a conclusion rests on
observational rather than experimental data. These habits, more than
any specific tool, are what should keep you in charge of your own
analysis as AI continues to change.

\section{Problems}

\begin{exercise}
Pick one bias from the catalog in Chapter~\ref{ch:data} (for example,
  survival bias or self-selection bias). Write a short prompt that you
  would give an AI assistant, describing a dataset and asking it to
  assess the risk of that specific bias. Then write, in your own
  words, what a good answer from the AI should include.
\end{exercise}

\begin{exercise}
Take an analysis from an earlier chapter in this book (for example,
  the logistic regression model in Chapter~\ref{ch:nyc}) and imagine
  asking an AI tool to critique it. List three questions you would
  want the AI to answer about the analysis, and explain why each
  question matters.
\end{exercise}

\begin{exercise}
Describe a realistic situation in which pasting a dataset into a
  general-purpose AI tool would raise a data-leakage or privacy
  concern. What alternative approach could you use to get help from
  AI without exposing the sensitive data itself?
\end{exercise}

\begin{exercise}
Suppose an AI tool recommends a t-test for a two-group comparison
  without being asked to justify the choice. Using what you learned in
  Chapter~\ref{ch:hypothesis}, write the follow-up questions you would
  ask the AI to confirm whether a t-test is really appropriate for
  your data.
\end{exercise}

\begin{exercise}
AI tools change quickly, but the principles in this chapter are meant
  to age well. Revisit this chapter's list of strengths and pitfalls a
  few months from now, after using an AI tool for a real analysis of
  your own. Which items still ring true? Would you add or remove
  anything based on your experience?
\end{exercise}


%%% Local Variables:
%%% mode: latex
%%% TeX-command-extra-options: "-shell-escape"
%%% TeX-engine: xetex
%%% TeX-master: "../sidsmain.tex"
%%% End:
